% Part: proof-theory
% Chapter: sequent-calculus
% Section: translations

\documentclass[../../../include/open-logic-section]{subfiles}

\begin{document}

\olfileid{pt}{seq}{tr}

\olsection{\Log{G1c} 与 \Log{G3c} 之间的转换}

我们之前提到，对于经典逻辑，\Log{G1c} 和 \Log{G3c} 都是可靠且完备的，
即，它们证明的恰好是在所有结构中皆为真的相继式。
特别地，它们证明相同的相继式。
现在我们无需绕道可靠性定理与完备性定理，纯证明论地即可证明这一事实。
这样的证明提供了将 \Log{G1c} 中的推导\emph{转换}为 \Log{G3c} 中的推导的方法，反之亦然。

\begin{prop}
若 $\Log{G1c} \Proves S$，则 $\Log{G3c} \Proves S$。
\end{prop}

\begin{proof}
  我们证明：若 $\pi$ 是 \Log{G1c} 中的一个推导，则存在 \Log{G3c} 中 $S$ 的一个推导~$\pi'$。
  我们对 $n = \pheight{\pi}$ 进行归纳。

  若 $n = 0$，则 $\pi$ 是一条公理。
  公理 $\lfalse \Sequent \quad$ 也是 \Log{G3c} 中的一条公理（取上下文 $\Gamma, \Delta$ 为空）。
  在 \Log{G3c} 中，允许公理 $!A \Sequent !A$ 中的 $!A$ 为任意公式，而不仅限于原子公式。
  但我们在 \olref[ex]{prop:G3c-axioms} 中已经证明了所有这样的相继式在 \Log{G3c} 中都是可推导的。
  
  若 $n > 0$，$\pi$ 以某个规则~$R$ 结束。
  归纳假设保证了该规则的前提在 \Log{G3c} 中是可推导的，即，直接子推导在 \Log{G3c} 中有相应的转换。
  如果规则~$R$ 也是 \Log{G3c} 中的规则，我们便将其应用于各前提转换后的推导。
  对于那些并非 \Log{G3c} 规则的规则，我们之前已经证明它们至少是可接纳的。
  因此，如果规则~$R$ 是 $\LeftR{\land}$ 或 $\RightR{\lor}$，则应用 \olref[adm]{prop:lor-G3-adm}；
  如果是弱化规则，则应用 \olref[adm]{prop:weak-G3c-adm}；
  如果是收缩规则，则应用 \olref[inv]{prop:G3c-cont-adm}。
\end{proof}

\begin{prop}
若 $\Log{G3c} \Proves S$，则 $\Log{G1c} \Proves S$。
\end{prop}

\begin{proof}
  我们再次对 $S$ 的一个推导~$\pi$ 的高度进行归纳。
  \Log{G3c} 的每条公理都可以在 \Log{G1c} 中由一条公理加上若干次 \LeftR{\Weakening} 和 \RightR{\Weakening} 规则的应用推导出来：
  \[
  \Axiom$!C \fCenter !C$
  \RightLabel{\LeftR{\Weakening}}
  \Deduce$!C, \Gamma \fCenter !C$
  \RightLabel{\RightR{\Weakening}}
  \Deduce$!C, \Gamma \fCenter \Delta, !C$
  \DisplayProof
\qquad
  \Axiom$\lfalse \fCenter$
  \RightLabel{\LeftR{\Weakening}}
  \Deduce$\lfalse, \Gamma \fCenter$
  \RightLabel{\RightR{\Weakening}}
  \Deduce$\lfalse, \Gamma \fCenter \Delta$
  \DisplayProof
  \]
  如果 $\pi$ 以某个规则~$R$ 结束，归纳假设给出了各前提在 \Log{G1c} 中的推导，我们可以将同样的规则~$R$ 应用于这些推导。
  例外是 \LeftR{\land} 和 \RightR{\lor} 规则，它们在 \Log{G1c} 与 \Log{G3c} 中有所不同。
  \Log{G3c} 版本的这些规则在 \Log{G1c} 中可按如下方式推导：
  \[
  \Axiom$!A, !B, \Gamma \fCenter \Delta$
  \RightLabel{\LeftR{\land}}
  \UnaryInf$!A \land !B, !B, \Gamma \fCenter \Delta$
  \RightLabel{\LeftR{\land}}
  \UnaryInf$!A \land !B, !A \land !B, \Gamma \fCenter \Delta$
  \RightLabel{\LeftR{\Contraction}}
  \UnaryInf$!A \land !B, \Gamma \fCenter \Delta$
    \DisplayProof
\qquad
  \Axiom$\Gamma \fCenter \Delta, !A, !B$
  \RightLabel{\RightR{\lor}}
  \UnaryInf$\Gamma \fCenter \Delta, !A \lor !B, !B$
  \RightLabel{\RightR{\lor}}
  \UnaryInf$\Gamma \fCenter \Delta, , !A \lor !B, !A \lor !B$
  \RightLabel{\RightR{\Contraction}}
  \UnaryInf$\Gamma \fCenter \Delta, !A \lor !B$
    \DisplayProof
  \]
\end{proof}

\end{document}